%!TeX root=../treasuretop.tex
\ArtChapter[The Cruise of the Coracle]{chap24head}

\lettrine[lines=5,image=true,findent=2pt]{chap24initI}{t} was broad day when I awoke and found myself tossing at the south-west end of Treasure Island. The sun was up but was still hid from me behind the great bulk of the Spy-glass, which on this side descended almost to the sea in formidable cliffs.

Haulbowline Head and Mizzen-mast Hill were at my elbow, the hill bare and dark, the head bound with cliffs forty or fifty feet high and fringed with great masses of fallen rock. I was scarce a quarter of a mile to seaward, and it was my first thought to paddle in and land.

That notion was soon given over. Among the fallen rocks the breakers spouted and bellowed; loud reverberations, heavy sprays flying and falling, succeeded one another from second to second; and I saw myself, if I ventured nearer, dashed to death upon the rough shore or spending my strength in vain to scale the beetling crags.

Nor was that all, for crawling together on flat tables of rock or letting themselves drop into the sea with loud reports I beheld huge slimy monsters—soft snails, as it were, of incredible bigness—two or three score of them together, making the rocks to echo with their barkings.

I have understood since that they were sea lions, and entirely harmless. But the look of them, added to the difficulty of the shore and the high running of the surf, was more than enough to disgust me of that landing-place. I felt willing rather to starve at sea than to confront such perils.

%   \begin{figure}[p]
%\centering
%\includegraphics[width=\linewidth]{chap24starveatsea}
%\caption[To confront such perils]{I felt willing rather to starve at sea than to confront such perils}
%\end{figure}   

\begin{bwbigpic}
	[\picsize]
	{chap24starveatsea}
	{I felt willing rather to starve at sea than to confront such perils}
	[To confront such perils]
\end{bwbigpic}

In the meantime I had a better chance, as I supposed, before me. North of Haulbowline Head, the land runs in a long way, leaving at low tide a long stretch of yellow sand. To the north of that, again, there comes another cape—Cape of the Woods, as it was marked upon the chart—buried in tall green pines, which descended to the margin of the sea.

I remembered what Silver had said about the current that sets northward along the whole west coast of Treasure Island, and seeing from my position that I was already under its influence, I preferred to leave Haulbowline Head behind me and reserve my strength for an attempt to land upon the kindlier-looking Cape of the Woods.

There was a great, smooth swell upon the sea. The wind blowing steady and gentle from the south, there was no contrariety between that and the current, and the billows rose and fell unbroken.

Had it been otherwise, I must long ago have perished; but as it was, it is surprising how easily and securely my little and light boat could ride. Often, as I still lay at the bottom and kept no more than an eye above the gunwale, I would see a big blue summit heaving close above me; yet the coracle would but bounce a little, dance as if on springs, and subside on the other side into the trough as lightly as a bird.

I began after a little to grow very bold and sat up to try my skill at paddling. But even a small change in the disposition of the weight will produce violent changes in the behaviour of a coracle. And I had hardly moved before the boat, giving up at once her gentle dancing movement, ran straight down a slope of water so steep that it made me giddy, and struck her nose, with a spout of spray, deep into the side of the next wave.

I was drenched and terrified, and fell instantly back into my old position, whereupon the coracle seemed to find her head again and led me as softly as before among the billows. It was plain she was not to be interfered with, and at that rate, since I could in no way influence her course, what hope had I left of reaching land?

I began to be horribly frightened, but I kept my head, for all that. First, moving with all care, I gradually baled out the coracle with my sea-cap; then, getting my eye once more above the gunwale, I set myself to study how it was she managed to slip so quietly through the rollers.

I found each wave, instead of the big, smooth glossy mountain it looks from shore or from a vessel's deck, was for all the world like any range of hills on dry land, full of peaks and smooth places and valleys. The coracle, left to herself, turning from side to side, threaded, so to speak, her way through these lower parts and avoided the steep slopes and higher, toppling summits of the wave.

»Well, now,« thought I to myself, »it is plain I must lie where I am and not disturb the balance; but it is plain also that I can put the paddle over the side and from time to time, in smooth places, give her a shove or two towards land.« No sooner thought upon than done. There I lay on my elbows in the most trying attitude, and every now and again gave a weak stroke or two to turn her head to shore.

It was very tiring and slow work, yet I did visibly gain ground; and as we drew near the Cape of the Woods, though I saw I must infallibly miss that point, I had still made some hundred yards of easting. I was, indeed, close in. I could see the cool green tree-tops swaying together in the breeze, and I felt sure I should make the next promontory without fail.

It was high time, for I now began to be tortured with thirst. The glow of the sun from above, its thousandfold reflection from the waves, the sea-water that fell and dried upon me, caking my very lips with salt, combined to make my throat burn and my brain ache. The sight of the trees so near at hand had almost made me sick with longing, but the current had soon carried me past the point, and as the next reach of sea opened out, I beheld a sight that changed the nature of my thoughts.

Right in front of me, not half a mile away, I beheld the \textit{Hispaniola} under sail. I made sure, of course, that I should be taken; but I was so distressed for want of water that I scarce knew whether to be glad or sorry at the thought, and long before I had come to a conclusion, surprise had taken entire possession of my mind and I could do nothing but stare and wonder.

The \textit{Hispaniola} was under her main-sail and two jibs, and the beautiful white canvas shone in the sun like snow or silver. When I first sighted her, all her sails were drawing; she was lying a course about north-west, and I presumed the men on board were going round the island on their way back to the anchorage. Presently she began to fetch more and more to the westward, so that I thought they had sighted me and were going about in chase. At last, however, she fell right into the wind's eye, was taken dead aback, and stood there awhile helpless, with her sails shivering.

»Clumsy fellows,« said I; »they must still be drunk as owls.« And I thought how Captain Smollett would have set them skipping.

Meanwhile the schooner gradually fell off and filled again upon another tack, sailed swiftly for a minute or so, and brought up once more dead in the wind's eye. Again and again was this repeated. To and fro, up and down, north, south, east, and west, the \textit{Hispaniola} sailed by swoops and dashes, and at each repetition ended as she had begun, with idly flapping canvas. It became plain to me that nobody was steering. And if so, where were the men? Either they were dead drunk or had deserted her, I thought, and perhaps if I could get on board I might return the vessel to her captain.

The current was bearing coracle and schooner southward at an equal rate. As for the latter's sailing, it was so wild and intermittent, and she hung each time so long in irons, that she certainly gained nothing, if she did not even lose. If only I dared to sit up and paddle, I made sure that I could overhaul her. The scheme had an air of adventure that inspired me, and the thought of the water breaker beside the fore companion doubled my growing courage.

Up I got, was welcomed almost instantly by another cloud of spray, but this time stuck to my purpose and set myself, with all my strength and caution, to paddle after the unsteered \textit{Hispaniola}. Once I shipped a sea so heavy that I had to stop and bail, with my heart fluttering like a bird, but gradually I got into the way of the thing and guided my coracle among the waves, with only now and then a blow upon her bows and a dash of foam in my face.

I was now gaining rapidly on the schooner; I could see the brass glisten on the tiller as it banged about, and still no soul appeared upon her decks. I could not choose but suppose she was deserted. If not, the men were lying drunk below, where I might batten them down, perhaps, and do what I chose with the ship.

For some time she had been doing the worse thing possible for me—standing still. She headed nearly due south, yawing, of course, all the time. Each time she fell off, her sails partly filled, and these brought her in a moment right to the wind again. I have said this was the worst thing possible for me, for helpless as she looked in this situation, with the canvas cracking like cannon and the blocks trundling and banging on the deck, she still continued to run away from me, not only with the speed of the current, but by the whole amount of her leeway, which was naturally great.

But now, at last, I had my chance. The breeze fell for some seconds, very low, and the current gradually turning her, the \textit{Hispaniola} revolved slowly round her centre and at last presented me her stern, with the cabin window still gaping open and the lamp over the table still burning on into the day. The main-sail hung drooped like a banner. She was stock-still but for the current.

For the last little while I had even lost, but now redoubling my efforts, I began once more to overhaul the chase.

I was not a hundred yards from her when the wind came again in a clap; she filled on the port tack and was off again, stooping and skimming like a swallow.

My first impulse was one of despair, but my second was towards joy. Round she came, till she was broadside on to me—round still till she had covered a half and then two thirds and then three quarters of the distance that separated us. I could see the waves boiling white under her forefoot. Immensely tall she looked to me from my low station in the coracle.

And then, of a sudden, I began to comprehend. I had scarce time to think—scarce time to act and save myself. I was on the summit of one swell when the schooner came stooping over the next. The bowsprit was over my head. I sprang to my feet and leaped, stamping the coracle under water. With one hand I caught the jib-boom, while my foot was lodged between the stay and the brace; and as I still clung there panting, a dull blow told me that the schooner had charged down upon and struck the coracle and that I was left without retreat on the \textit{Hispaniola}.